\providecommand{\cH}{\mathcal{H}}
\providecommand{\cC}{\mathcal{C}}
\providecommand{\Real}{\mathrm{Re}}
\providecommand{\Imag}{\mathrm{Im}}
\providecommand{\Pfaff}{\mathrm{Pf}}
\providecommand{\DeltaFive}{\Delta{}_{5}}
\providecommand{\PhiTen}{\Phi{}_{10}}
\providecommand{\CoHA}{\operatorname{CoHA}}

\section*{The $\wedge^2$ Transfer and the $\DeltaFive$ Reflection System}

\subsection*{The Pfaffian Lattice}

\noindent\textbf{Theorem (Pfaffian transfer).}
Let $\Lambda^4 = \bigoplus_{i=1}^{4}\mathbb{Z}e_i$. The Pfaffian pairing
\[
u \wedge v = (u,v)\, e_1\wedge e_2\wedge e_3\wedge e_4 \in \wedge^4\Lambda^4
\]
is $\mathrm{SL}_4(\mathbb{Z})$-invariant, integral, symmetric, even, unimodular,
of signature $(3,3)$ on $\wedge^2\Lambda^4$. Fix
\[
q = e_1\wedge e_3 + e_2\wedge e_4 .
\]
For $x,y\in\Lambda^4$, the scalar
\[
B_q(x,y)e_1\wedge e_2\wedge e_3\wedge e_4 = -x\wedge y\wedge q
\]
defines the integral symplectic form preserved by $\mathrm{Sp}_4(\mathbb{Z})$.
The stabiliser of $q$ inside $\mathrm{SL}_4(\mathbb{Z})$ is therefore
\[
\{g \in \mathrm{SL}_4(\mathbb{Z})\,:\, (g\wedge g)(q) = q\}
\simeq \mathrm{Sp}_4(\mathbb{Z}).
\]

\subsection*{The Orthogonal Complement}

The Gritsenko--Nikulin lattice is the orthogonal complement
\[
\Lambda^{3,2} \;=\; (e_1\wedge e_3 + e_2\wedge e_4)^{\perp}
\;\simeq\; \Lambda^{1,1}\oplus\Lambda^{1,1}\oplus[2],
\]
of signature $(3,2)$. A convenient basis is
\[
f_1 = e_1\wedge e_2,\quad f_2 = e_2\wedge e_3,\quad
f_3 = e_1\wedge e_3 - e_2\wedge e_4,\quad
f_{-2} = e_4\wedge e_1,\quad f_{-1} = e_4\wedge e_3.
\]
The rank-$3$ hyperbolic sublattice
\[
\Lambda^{2,1}=\Lambda^{1,1}\oplus[2]
\simeq\mathbb{Z}f_2\oplus\mathbb{Z}f_3\oplus\mathbb{Z}f_{-2}
\]
is the reflection lattice on which the real simple roots of
$\mathfrak{g}_{\DeltaFive}$ live.

\subsection*{The Type-IV Domain}

The real orthogonal group $\mathrm{O}_{\mathbb{R}}(\Lambda^{3,2})$
acts on
\[
\cH^{\mathrm{IV}} = \bigl\{Z\in\mathbb{P}(\Lambda^{3,2}\otimes\mathbb{C})
\,:\, (Z,Z)=0,\ (Z,\overline{Z})<0\bigr\}
= \cH^{\mathrm{IV}}_{+}\cup\cH^{\mathrm{IV}}_{-}.
\]
In the $(f_i)$-basis a point of $\cH^{\mathrm{IV}}_{+}$ is
$Z = ((z_2^2 - z_1 z_3),\, z_3,\, z_2,\, z_1,\, 1)\cdot z_0$ with
$\Imag(z_1)>0$. The isotropy relation $(Z,Z)=0$ is automatic; the
negativity $(Z,\overline{Z})<0$ reduces to
\[
y_1 y_3 - y_2^{\,2} > 0, \qquad y_i = \Imag z_i,
\]
with $\Imag\left(\begin{smallmatrix}z_1&z_2\\z_2&z_3\end{smallmatrix}\right)$
positive definite. Thus $\cH^{\mathrm{IV}}_{+}$ is the type-IV
realisation of the Siegel upper half space $\mathbb{H}_2$.
The index-$2$ subgroup $\mathrm{O}_{\mathbb{R}}(\Lambda^{3,2})_+$ is the
one fixing $\cH^{\mathrm{IV}}_+$; its kernel on $\cH^{\mathrm{IV}}_\pm$
is $\pm I_5$, and since $\Lambda^{3,2}$ is odd-dimensional,
$\mathrm{O}_{\mathbb{R}}(\Lambda^{3,2})_+ = \pm I_5\cdot
\mathrm{SO}_{\mathbb{R}}(\Lambda^{3,2})_+$.

\subsection*{The Integral Transfer}

The exterior-square representation restricts to the integral
orthogonal action
\[
\wedge^2\,:\, \mathrm{Sp}_4(\mathbb{Z})/\{\pm I_4\}
\;\xrightarrow{\ \sim\ }\;
\mathrm{O}(\Lambda^{3,2})_+/\{\pm I_5\}.
\]
This transfer identifies the $\mathrm{Sp}_4(\mathbb{Z})$ modular form
\(\DeltaFive\) with its reflective orthogonal avatar on
$\cH^{\mathrm{IV}}_+$. The Borcherds product and the
Weyl--Kac--Borcherds sum therefore describe the same automorphic
denominator.

\subsection*{Reflections and the Character of $\DeltaFive$}

\noindent\textbf{Proposition (reflection character).}
For $\alpha\in\Lambda^{2,1}$ primitive with
$(\alpha,\alpha)>0$ and $(\alpha,\alpha)\mid 2(\Lambda^{2,1},\alpha)$
the reflection is
\[
s_{\alpha}(x) = x - \frac{2(x,\alpha)}{(\alpha,\alpha)}\alpha, \qquad
\alpha^2 = 2 \;\Rightarrow\; s_{\alpha}(x) = x - (x,\alpha)\alpha.
\]
Maass' multiplier formula splits the square-$2$ roots into two
characters for $\DeltaFive$:
\begin{itemize}
\item $\alpha\in\{\delta_1 = 2f_2 - f_3,\ \delta_2 = 2f_{-2} - f_3,\
\delta_3 = f_3\}$: \quad $\DeltaFive(s_\alpha z) = -\DeltaFive(z)$,
\item $\alpha\in\{f_{-2}-f_2,\ f_2-f_3,\ f_2+f_3\}$: \quad
$\DeltaFive(s_\alpha z) = +\DeltaFive(z)$.
\end{itemize}
The three anti-invariant generators $\{\delta_1,\delta_2,\delta_3\}$
are the real simple roots of the generalized Borcherds--Kac--Moody
superalgebra $\mathfrak{g}_{\DeltaFive}$;
their Gram matrix is $2$ on the diagonal and $-2$ off-diagonal, the
hyperbolic Cartan matrix that controls $W^{(2)}_{\mathrm{II}}$.
The invariant reflections generate $W^{(2)}_{\mathrm{I}}$; the
anti-invariant reflections generate the Weyl group entering the
denominator identity. Concretely, for $U\in\mathrm{GL}_2(\mathbb{Z})$,
the reflection lifts through
$\wedge^2\left(\begin{smallmatrix}0&U\\{}^tU^{-1}&0\end{smallmatrix}\right)$,
and the multiplier gives the displayed sign.

\subsection*{The Denominator and the BKM Weight}

\noindent\textbf{Theorem (primitive denominator).}
The three roots $\delta_i$ generate the Weyl vector
\[
\rho = \tfrac{1}{2}(\delta_1+\delta_2+\delta_3)
= f_2 - \tfrac{1}{2}f_3 + f_{-2}, \qquad
(\rho,\delta_i)=-1 .
\]
The Borcherds denominator identity is
\[
\DeltaFive(z) \;=\; e^{-2\pi i(\rho,z)}
\prod_{\alpha\in\Delta_+}\bigl(1 - e^{-2\pi i(\alpha,z)}\bigr)^{\mathrm{mult}\,\alpha}
\]
with root multiplicities read from the Fourier coefficients of
$\phi_{0,1}$. Borcherds' weight theorem gives
\[
\kappa_{\mathrm{BKM}}(\Phi_1)
=\kappa_{\mathrm{BKM}}(\DeltaFive)
=\frac{c_1(0)}{2}
=\frac{10}{2}
=5 .
\]
The exterior-square transfer above identifies the Siegel form and the
orthogonal reflective form before the denominator is read, so the
weight-$5$ Borcherds invariant and the reflection character refer to
one modular form.

The square scalar belongs to the doubled Jacobi input:
\[
\mathrm{BorLift}\bigl(2\,\phi_{0,1}\bigr)=\PhiTen^{\mathrm{un}}
=\DeltaFive^{\,2}, \qquad
\mathrm{BorLift}^{\mathrm{prim}}\bigl(\phi_{0,1}\bigr)=\DeltaFive .
\]
Thus the primitive signed root characters of
$\mathfrak{g}_{\DeltaFive}$ are read from $\phi_{0,1}$, while the
Igusa scalar trace reads the square.

\noindent\textbf{Corollary (primitive CHL constants).}
For the primitive CHL denominator ladder
$N\in\{1,2,3,4,6\}$,
\[
\bigl(c_N(0)\bigr)_{N=1,2,3,4,6}=(10,8,6,4,2),\qquad
\bigl(\kappa_{\mathrm{BKM}}(\Phi_N)\bigr)_{N=1,2,3,4,6}=(5,4,3,2,1).
\]
These constants use the CHL/Borcherds denominator normalisation.
The Gritsenko--Clery catalogue is an eight-row atlas with its own row
normalisations and cover groups; a comparison with the CHL ladder
requires an explicit row map and a scalar normalisation.

\subsection*{Chiral Scope}

\noindent\textbf{Proposition (native level).}
The $K3\times E$ Hall--Drinfeld/Borcherds branch is built from the
positive Hall half, its completed double, and the reflective denominator
above. The self-mirror K3 Yangian branch has a distinct Mukai-current
input. For a CY input of dimension $d\geq 3$, the specialised output
$A=\Phi_d(\mathcal{C})$ is an $E_1$-chiral algebra on the chosen curve;
$E_2$ structure enters through centre, Drinfeld-double, module, or
stable-envelope enhancement layers.

The affine local model obeys
\[
\CoHA(\mathbb{C}^{3})=Y^+(\widehat{\mathfrak{gl}}_1).
\]
The route to $\mathcal{W}_{1+\infty}[\lambda]$ factors through the
Drinfeld double, centre, completion, and evaluation map. The compact
and specialised values attached to $K3\times E$ are distinct
constructions:
\[
\begin{array}{c|c|l}
\text{invariant} & \text{value} & \text{construction}\\
\hline
\kappa_{\mathrm{cat}}(K3\times E) & 0 &
\chi(\mathcal{O}_{K3})\chi(\mathcal{O}_{E})\\
\kappa_{\mathrm{ch}}(K3\times E) & 0 &
\sum_q(-1)^q h^{0,q}(K3\times E)\\
\kappa_{\mathrm{ch}}^{\mathrm{Heis}}(K3\times E) & 3 &
\text{Heisenberg--Mukai specialisation}\\
\kappa_{\mathrm{BKM}}(\DeltaFive) & 5 &
c_1(0)/2\\
\kappa_{\mathrm{fiber}}(K3) & 24 &
\mathrm{rk}\,\widetilde H(K3,\mathbb{Z})
\end{array}
\]
The categorical Euler is K\"unneth-multiplicative:
\[
\kappa_{\mathrm{cat}}(K3\times E)
=\chi(\mathcal{O}_{K3})\chi(\mathcal{O}_{E})
=2\cdot 0
=0 .
\]
The compact chiral Hodge supertrace is
\[
\kappa_{\mathrm{ch}}(K3\times E)
=h^{0,0}-h^{0,1}+h^{0,2}-h^{0,3}
=1-1+1-1
=0 .
\]
The Heisenberg value $3$ belongs to the Mukai-enhanced specialisation,
the value $5$ belongs to the Borcherds denominator, and the value $24$
belongs to the K3 fibre lattice. Repetition of one $\Phi_3$ input
cannot produce this table.

\subsection*{Proof Obligations}

The compact Hall comparison requires the following constructions.
\begin{enumerate}
\item Construct an oriented compact critical CoHA positive half for
$K3\times E$ with the reduced obstruction theory and Serre-equivariant
gluing data.
\item Produce a non-degenerate Hopf pairing and completion for the
Hall half, yielding the Hall--Drinfeld double carrying the universal
$R$-matrix.
\item Identify the primitive Hall signed character with the
Borcherds denominator character of $\mathfrak{g}_{\DeltaFive}$,
including the parity fixture that turns signed superdimension into
ordinary dimension where needed.
\item Supply any comparison between a CHL primitive denominator and a
Gritsenko--Clery catalogue row by an explicit row map, multiplier
system, and scalar normalisation.
\end{enumerate}
The Pfaffian transfer, the reflection character, and the numerical
constants above are denominator-level statements and are independent
of these compact Hall constructions.
